\documentclass{article}
\title{統計で読むニュースの嘘 — 数字の出し方を見抜く}
\author{TeX 図書館}

\begin{document}

\section{数字は誰が・どう選ぶか}

ニュースの見出しには数字が並ぶ。「7割」「急増」「5倍」。数字は嘘をつかないと言われる。しかし数字は、事実の候補から誰かが1つ選んだものである。選ぶ作業は記事の外で起きる。発表資料、調査報告、社内集計。どれを抜粋するかは書き手の判断だ。

判断には都合が働く。売りたい企業、支持率を上げたい政治家、速報を競う新聞。同じ集計からでも、都合の良い数字は簡単に見つかる。この本は、その\textbf{選び方と見せ方}を外から疑うための本だ。

公式は一つも置かない。必要なのは足し算と割り算、そして疑う習慣だけである。各章の末尾には【演習】を置く。答えを見せる前に、実際の記事を1つ思い出してから解答の指針を開くと効果が上がる。最後の章には、そのまま使えるチェックリストと報道機関への問い合わせの文面を収めた。

都合の良い数字の選び方には、いくつかの型がある。どれも統計の誤りではない。事実として立つ数字だけれど、選び方が一方に傾いている。

\begin{itemize}
\item \textbf{指標のすり替え}: 増えている指標だけを報告する。減っている指標は口に出さない。
\item \textbf{カテゴリの細分化}: カテゴリを細かく割れば、どこかで「第1位」が作れる。
\item \textbf{期間の選択}: 始点と終点を動かせば、傾向は自由に描ける(第2章)。
\end{itemize}

\begin{quote}
\textbf{【例】}\ 健康グッズの広告が「猫背改善ブランディング第1位」と謳っていた。調べると、主催者が作ったカテゴリに参加企業は2社だけだった。数字は偽っていない。ただ、第1位という語だけが記事に残る。
\end{quote}

自治体の例も似ている。「ゴミ収集1戸あたりのコストを1割下げた」と発表した市がある。人口減少で収集量自体が減っていたため、総額はむしろ増えていた。指標を1つ選べば、同じ仕事の評価は正反対にもなる。

\begin{quote}
\textbf{【注意】}\ 報道された数字が計算間違いであることは、実はあまりない。多いのは、正しい数字から都合の良い切り口を選んでいる場合だ。だから最初に問うのは「合っているか」ではなく「なぜこの切り口か」である。
\end{quote}

この本では、記事を読むたびに次の3問を唱える習慣を作る。

\begin{enumerate}
\item \textbf{誰が}この数字を選んだか。発表者と利害の関係は何か。
\item \textbf{何と}比べているか。比較対象は自分で選んだものか。
\item \textbf{母数は}何か。割合の分母は誰で、何人か。
\end{enumerate}

3問はそれぞれ後の章に受け渡される。第2章と第3章で比較、第3章で母数、第4章以降で指標の中身を掘る。総合章では、この3問を12個のチェック項目に展開する。

この本の使い方は次のとおりである。

\begin{itemize}
\item 各章の【演習】は、手元にある記事1本に向かって答えると効く。
\item 数字を出す側に回ったときも、同じ3問を自分に向ける。
\item 3問のどれも埋まらない記事は、結論を保留して読んでよい。
\end{itemize}

\begin{quote}
\textbf{【定義】}\ \textbf{都合の良い数字}とは、嘘はついていないが、指標・比較対象・期間を利害に沿って選んだ数字を指す。事実は保たれ、印象だけが動く。この点が、明確な虚偽より読み抜きにくい。
\end{quote}

\section{比較のトリック}

比較なしに数字は語れない。「増加した」は必ず何かと比べた増加だ。問題は、比較対象を書き手が選べる点にある。

\subsection{比較対象のすり替え}

新製品の発表で「従来モデルの3倍速」と書かれていたら、従来モデルが何かを確認する。5年前の廉価版と比べれば、倍率はいくらでも膨らむ。比較対象を最安の旧機種に置くのが定石になっている。

\begin{quote}
\textbf{【例】}\ 扇風機の広告に「従来品より静音」とあった。従来品は前世代の最下位モデルで、主力機どうしでは差が半分しかない。正しい記述でも、比較対象の選択だけで印象は二転する。
\end{quote}

企業どうしの比較も同じ。「主要3社中最静か」と謳うとき、3社の顔ぶれは誰が決めたか。比較対象を足せば順位は下がり、外せば順位は上がる。

\begin{quote}
\textbf{【定義】}\ \textbf{比較対象のすり替え}とは、データを一切動かさずに、比較の相手を選ぶだけで倍率や順位を動かすことを指す。数字はすべて正しく、結論だけが傾く。誤りとして訂正を求めることができないのが厄介な点だ。
\end{quote}

\subsection{期間の切り取り}

「売上は過去最高を更新」の始点を問う。10年前の底値から数えれば、停滞している事業も右肩上がりに見える。始点と終点の両方が必要だ。片方だけの期間指定は、通常どの記事でも省略される。「いつから」を先に書き留めておくと、次の記事との比較が速い。

\begin{quote}
\textbf{【注意】}\ 期間は数字と一緒には書かれない。「昨年来」「過去最高」という語は、実際にはどの区間かを本文で確かめたい。確認しないまま受け取ると、書き手が選んだ期間を自分の記憶と同じものだと思い込む。
\end{quote}

\subsection{前年比と対前月比}

同じ出来事でも、比較の基準を変えると見出しは変わる。

\begin{itemize}
\item \textbf{前年比}: 同じ季節と比べるため季節の山を消せる。ただし前年が異常だった年は、その影響を背負う。
\item \textbf{対前月比}: 直近の動きに敏感。季節の山と谷をそのまま拾う(12月は11月から必ず下がる。4月は新年度で跳ねる)。
\end{itemize}

前年比を見たら、前年の水準も聞きたい。前年が特別に落ち込んでいた年は、普通に戻っただけで大きな伸びになる。季節で動く分野ほど、この罠は深い。旅行、電力、農産物が典型である。

どちらを選ぶかは物語次第だ。好調を強調したいなら前年比、警戒を煽りたいなら対前月比が選ばれやすい。

\begin{quote}
\textbf{【例】}\ 同じ売上について「前年同月比で2割増」と「直近1か月では3\%減」の2つの見出しが出た。数字はどちらも正しく、どちらも本文に書いてある。見出しだけで受け取れば、読者は別々の会社を想像する。
\end{quote}

\begin{quote}
\textbf{【演習】}\ 記事が「対前年同月比プラス40\%」だけを大きく書き、対前月比は本文の隅に小さく載せていた。前年と前月で、それぞれ何が起きていたと推測されるか。2つ挙げよ。
\end{quote}

\textbf{解答の指針:} 前年は特別に落ち込んでいた可能性(工場の停止、輸入の停止、前の年の大量発注の反動)。前月からは減っているので、直近の勢いは落ちている可能性。前年比の大きな伸びは、前年の水準が低かっただけかもしれない。

\section{母数と構成比}

「7割が支持」の次に聞くべきは、母数である。誰の7割か。

\begin{quote}
\textbf{【定義】}\ \textbf{母数}とは、割合の分母にあたる人数や件数を指す。支持率なら回答者数、シェアなら市場全体の規模。「7割」という比率だけでは、280人と28人のどちらなのか分からない。
\end{quote}

経営者400人に聞いた支持率7割は、400人中280人の意思である。総人口の7割とは全く別の話になる。街頭インタビュー100人の7割は、答えてくれた100人の7割だ。断った人や無視した人の分は母数から消えている。

構成比でも同じことが起きる。「支持7割、反対2割、残り1割」の残りが無回答なら、無回答を分母から外して集計し直すと支持率は上がる。分母の定義は集計者が決められる。

構成比の合計も手がかりになる。合計が100から大きく外れているときは、丸め方の問題か、選択肢の重複か、抜き取りが起きている。小さな食い違いでも、記事の作り方を疑う入口にはなる。

\begin{quote}
\textbf{【注意】}\ 街頭インタビューと、自発的参加のウェブ調査は偏りやすい。意見が強い人ほど答えるからだ。答えた人が母集団を代表しているという保証は、選抜方法を書かない限りない。
\end{quote}

サンプル数も必ず聞く。30人中21人が支持なら7割だが、次の調査では21人中13人で6割に落ちる。少数のサンプルは、割合が10ポイント単位で揺れる。母数が小さい7割は、ほとんど何も約束しない。

\begin{quote}
\textbf{【例】}\ 「国内シェア7割」の母数が、市場全体の台数なのか、対象カテゴリに限った台数なのかで、意味は変わる。カテゴリを狭めた瞬間にシェアは跳ね上がる。第1章の細分化と、対で読むとよく分かる。
\end{quote}

記事に向ける質問は次の順になる。

\begin{enumerate}
\item 分母は誰か(回答者か、全員か)。
\item 実数は何人か(割合とセットで実数が書かれているか)。
\item 何人から選ばれたか(無作為か、自発参加か)。
\item 無回答と未定はどこに入ったか。
\end{enumerate}

\section{平均の罠}

平均と中央値は別の道具である。平均は全部の値を合計して人数で割る。中央値は値を並べて真ん中を取る。データが左右に歪んでいるとき、両者は離れる。

\begin{quote}
\textbf{【例】}\ 新車の平均購入価格が400万円から450万円に上がったと報じられた。「新車が値上げ」という見出しが続いた。ところが中央値は250万円のままだった。買った人の層が高級車に寄っただけで、同じ車は値上がりしていない。
\end{quote}

平均は全データの値を1つずつ使う。だから1つの極端な値に引っ張られる。中央値は端の値がいくらかを問わない。故障1件のような出来事には動かない。

\begin{quote}
\textbf{【例】}\ 平均待ち時間4分の店で、機材故障により1件だけ26分かかった日がある。その日の平均は跳ね上がるが、並んだ人の半数は相変わらず4分前後に収まっている。平均が語るのは店の日常ではない。
\end{quote}

所得、家賃、待ち時間、売上、利用時間。これらは右に長い尻尾を引くデータだ。だから平均だけが書かれた報道は、分布を隠している疑いを持つ価値がある。中央値と実数が添えられている報道のほうが忠実である。

\begin{quote}
\textbf{【例】}\ あるアプリの「1日平均利用者10万人」という数字には、土日の30万人と平日の4万人が入っていた。平均はどの曜日にも当てはまらない。広告の見積もりは実態より大きく見える。
\end{quote}

平均を基準にした評価にも注意が要る。「平均より下の人が半数以上いる」と報じられても、不思議ではない。平均は合計を人数で割った数字だから、下にいる人が半分を超えることは珍しくない。平均との差は、優劣の判定には使えない。

\begin{quote}
\textbf{【注意】}\ 平均を下回ったというだけで失敗と言えないのと同じで、平均を超えても上位とは限らない。外れ値1つが平均を大きく動かすので、平均との距離は実感と合わないことが多い。中央値と聞いてから判断したい。
\end{quote}

記事に向ける質問は3つである。平均か中央値か。何件のデータか。最大と最小はどこか。

\begin{quote}
\textbf{【演習】}\ ある店のレジ待ち時間が5件で 3分、3分、4分、4分、26分だった。平均と中央値を言葉で説明し、どちらがその店の日常を表すか答えよ。
\end{quote}

\textbf{解答の指針:} 5件の合計は40分で、それを5人で割ると平均は8分である。値を並べた中央は4分だ。故障の1件だけが平均を2倍に押し上げているので、日常を表すのは中央値の4分である。

\section{相対リスクと絶対リスク}

「発生リスクが5倍」は\textbf{相対リスク}である。元の発生率を知らない限り、この数字では判断できない。

\begin{quote}
\textbf{【定義】}\ \textbf{相対リスク}は2つの群の発生率の\textbf{比}、\textbf{絶対リスク}はそれぞれの\textbf{発生率そのもの}を指す。比は大きく見えやすく、見出し向きだ。実際の重さを表すのは発生率である。
\end{quote}

\begin{quote}
\textbf{【例】}\ 特定の飲料を飲む群1万人で20件、飲まない群1万人で4件のがん発症があった。20は4の5倍だから「5倍の危険」と書ける。しかし発生率は0.2\%と0.04\%で、差は0.16\%である。2万人のうち16人に1人の差だ。
\end{quote}

比だけが大きく見えるのは、元の発生率が小さいときほど起きやすい。まれな事象の小さな変化は、比にすると急に大きく、差にするとほとんど見えなくなる。

増加の報道でも同じことが起きる。「件数が3倍に増加」の裏に、1件から3件という実数が隠れていることがある。比が大きく見えるほど、実数は小さく書かれやすい。本文の最後に置かれた実数を、必ず拾いに行く。

\begin{quote}
\textbf{【注意】}\ 比が同じでも、期間が違えば意味は変わる。1年あたりの差と、5年間の合計の差は別の話である。記事が期間を書かないときこそ、比だけを信じてはいけない。
\end{quote}

発生率の言葉は翻訳して読みたい。「10万人に20人」は0.02\%、「5000人に1人」も同じ意味である。自明でない単位(1日あたり、1年あたり、服用1回あたり)が書かれているかを確かめるのも忘れずに。

\begin{itemize}
\item \textbf{見出しは相対}: 「3倍に増加」「リスクが上がる」は比の話。
\item \textbf{本文は絶対}: 「何件あったか」「何人のうち何人か」は発生率の話。
\item \textbf{期間の長さ}: 10年続けての差と、1回の差を同じ文章で語られていないか。
\end{itemize}

この3つが揃っていない記事は、比の大きさだけで重さを伝えている可能性が高い。

\begin{quote}
\textbf{【演習】}\ あるサプリの調査で、1万人中30人が体調不良を訴え、プラシーボ群では1万人中10人だった。相対と絶対をそれぞれ言葉で書き、見出しにできそうな2種類を挙げよ。
\end{quote}

\textbf{解答の指針:} 相対は30の10の3倍、つまり3倍である。絶対は0.3\%と0.1\%で、差は0.2\%だ。見出しは「体調不良の訴え3倍に」(相対)と「2万人に1人の差」(絶対)のどちらでも作れるが、与える印象は正反対になる。

\section{相関と因果}

「Aが多い地域ではBも多い」という報告をよく見る。この文から言えるのは、AとBが一緒に動く関係がある、だけである。

言えないことは3つある。AがBを起こしたこと。BがAを起こしたこと。そして、この関係が偶然ではないこと。

\subsection{時間の向き}

\begin{quote}
\textbf{【例】}\ オンラインゲームのプレイ時間が長い生徒ほど成績が低い、という調査が載った。原因の向きは2通り考えられる。ゲームに逃げるから成績が落ちるのか、睡眠時間を削ってゲームと成績が両方落ちるのか。観察だけでは向きが決まらない。
\end{quote}

\textbf{逆因果}はこの向きの問題である。関係は本物でも、原因と結果が逆かもしれないと考える。

\subsection{第3の要因}

\textbf{交絡}は第3の要因の問題だ。睡眠時間、家庭環境、学習意欲。これらが両方を動かすと、AとBは関係なく並んで見える。調査に載っていない要因ほど、記事本文にも登場しない。

\subsection{集計の単位}

\begin{quote}
\textbf{【例】}\ 「コンビニの多い街では平均寿命が短い」という見出しがあった。店は需要で選ばれる。夜に食べる人が多い街ほど店が集まる。店が寿命を縮めているのではない。順番が逆の疑いと、街の構成要素の両方を問うたい。
\end{quote}

地域の数字は個人の数字でもある。「歩く人の多い地域では糖尿病が少ない」を、個人の予防効果と読み替えてはいけない。地域単位の集計には、高齢者比率、所得、医療機関の数が混ざっている。集計の単位が変われば、見え方も変わる。

\begin{quote}
\textbf{【注意】}\ 相関の強さは因果の強さではない。観察データだけでは「効いた」は言えない。効きを見るには、条件を揃えて比べる仕組み(人を無作為に分けて比較するなど)が要る。その手続きを書いた報道は少ない。
\end{quote}

記事に向ける質問は3つである。時間がどうなっているか(先にAか、先にBか)。両方を動かす隠れた要因は何か。そして、比較する仕組みがあったか。

\section{調査の読み方}

世論調査の数字は、設問の作られ方で動く。同じテーマでも、聞き方を変えると答えは変わる。

\subsection{設問順と回答の余地}

候補者の名前を先に並べた設問では、前に出た名前が有利になる。「他にもありますか」と聞くか「以下から選んでください」と聞くかで、名前の出方は変わる。

4択で「どちらとも言えない」を含めると、支持率は下がる。支持と反対の2択に絞ると、迷った人までどちらかに押し込まれる。同じ55\%でも、設問が違えば水準は違う。

\begin{quote}
\textbf{【例】}\ 同じ候補者3人への支持率を、名前の並び順を変えて2種類の調査を作ると、先に挙がった人が数ポイント有利になった。公表された調査票には、質問の順序まで書いてある。
\end{quote}

\begin{quote}
\textbf{【注意】}\ 設問文の文言も違う。「支持するか」か「好感を持っているか」では、答えの温度が変わる。記事に設問文が載っていないときこそ、出典を遻って確かめたい。
\end{quote}

\subsection{誰が答えたか}

調査の質は、母数だけでは決まらない。何人に声をかけ、何人から答えてもらったか。ここも問いたい。回答の割合が低いほど、答えた人は偏る。応募者だけのアンケートは、意見が強い人が集まりやすい。

\subsection{誤差と煽り方}

調査は一部の人から集める。だから必ずぶれる。1000人規模の調査なら、数ポイントぶれるのが普通である。52\%対48\%の差は、次の調査で逆転してよい水準だ。

選挙報道の「逆転」「急伸」は、このぶれの範囲で書かれがちである。期日前の票と当日の票で有権者の構成も変わる。調査の実施時期と投票日までの間の出来事も、数字に混ざる。

調査会社名と実施日は、記事の隅に小さく書かれていることが多い。この2つは最低限拾ってから、数字の意味を判断したい。

\begin{quote}
\textbf{【例】}\ 「支持率で逆転」という見出しの裏を辿ると、前回調査との差は2ポイントで、調査対象は800人だった。見出しは誤差の内側で作られている。本文には調査会社名と実施日が小さく書いてある。
\end{quote}

\begin{quote}
\textbf{【演習】}\ 調査Aは1000人中で52\%対48\%、調査Bは300人中で55\%対45\%だった。どちらの数字を厚く見るべきか。理由を2つ書きよ。
\end{quote}

\textbf{解答の指針:} 母数の大きい調査Aのほうがぶれが小さい。Bは300人なので、数ポイントの差は当たり前の範囲に収まる。もう1つの理由は、両調査の設問文と選抜方法が揃っているかを確かめる必要があることだ。母数だけで優劣は決まらないが、まず疑うのはAよりBである。

\section{「統計的に有意」の誤読}

報道に「統計的に有意な差」という語が登場したら、次の定義を思い出してほしい。

\begin{quote}
\textbf{【定義】}\ \textbf{統計的に有意}とは、差がないと仮定した場合に、観測ほどの差が偶然で出る確率が、事前に決めた基準より小さい状態を指す。基準は慣習的に0.05、つまり20回に1回以下に置かれることが多い。これは差の\textbf{大きさ}を述べているのではない。
\end{quote}

直感的にはこうだ。「差が起きない世界で、これ以上に大きな差が20回に1回も出ないなら、差なしという説明は乏しい」。基準は約束事ではなく、あくまで目安である。0.049と0.051に実質的な違いはないのに、語法だけが変わる。

基準の0.05は慣習にすぎない。0.01で示す分野もあれば、0.10で示す分野もある。事前に基準を決めていたか、データを見てから決めていないか。この違いは記事にほぼ書かれないため、聞きたい項目である。

\subsection{有意は重要と同義ではない}

10万人の調査では、ごく小さな差でも有意になる。母数が大きければ、わずかなズレも偶然では説明しづらくなる。逆に100人ほどの調査では、目に見える差でも有意にならないことがある。証拠が足りないだけである。

\begin{quote}
\textbf{【例】}\ 10万人を対象にした調査で、2群の差が0.5ポイントでも「有意な差」と報告された。一方、150人の小規模調査では、目に見える差が「有意でない」とされた。差の判定と、差の評価は別の話である。
\end{quote}

\begin{quote}
\textbf{【注意】}\ 「有意だった」は「効果が大きい」の意味ではない。「有意でなかった」は「効果がない」の意味でもない。差の大きさと、検査に使った人数は、記事に書いていなければ自分で聞きたい項目である。
\end{quote}

\subsection{再現性を疑う}

初の1報は慎重に読む。同じデータを何度も切り分けて、良い結果が出た切り口だけを報告することがありえる。1つの研究だけが良質な結果を出すとき、次の研究が同じ結果を出すか確かめたい。

\begin{itemize}
\item 何人のデータか。少人数で出た一発の結果か。
\item 何回調べたか。調べる回数が多いほど、当たりは出る。
\item 別の研究も同じ向きを示しているか。
\item 差の大きさは、実務で意味を持つ水準か。
\end{itemize}

\section{記事を検証する}

前8章の質問を、読む場面で使える形に並べる。記事を1本読むたびに、上から順に確かめる。

\begin{enumerate}
\item 比較対象は誰が選んだか。
\item 期間の始点と終点はどこか。
\item 前年比か、対前月比か。どちらが都合が良いか。
\item 母数は何人か。実数が書かれているか。
\item 何人から、どう選ばれた人か。
\item 無回答と未定はどこに入ったか。
\item 平均か中央値か。極端な値が混ざっていないか。
\item 相対か絶対か。発生率そのものは何か。
\item 比較の期間はどれ分か。
\item 相関か因果か。時間の向きが書かれているか。
\item 誤差は何ポイントか。差はその範囲の外か。
\item 誰に都合が良い発表か。利害は何か。
\end{enumerate}

12項目は3問の延長線上にある。比較関係が6項目、母数と選抜が2項目、指標の中身が3項目、発表者の利害が1項目だ。

\subsection{問い合わせの型}

確認したい項目が埋まらないときは、出典に当たるのが早い。報道機関と調査元は、通常この問い合わせに答える。

\begin{quote}
\textbf{【例】}\ 「お世話になっております。○月○日の記事『△△』について、統計の出典を確認させてください。次をお知らせいただけますでしょうか。調査主体と実施日、調査対象者数と選抜方法、設問文の全文、無回答と未定の扱い、集計値が平均か中央値か。回答が難しい場合は、参照した一次資料の名称だけでも結構です。よろしくお願いいたします。」
\end{quote}

一次資料に辿り着けば、この本で見てきた穴はほとんど埋まる。資料が公開されていないとき、それ自体が回答である。読者は結論を保留してよい。

問い合わせの要点は3つである。

\begin{itemize}
\item 質問は項目名で箇条書きにする。曖昧な依頼には、曖昧な回答しか返ってこない。
\item 一次資料の名称をあわせて求めれば、公開されていれば自分で確かめられる。
\item 回答がない場合は、非公開という情報を得たと捉える。
\end{itemize}

\begin{quote}
\textbf{【総合演習】}\ 見出し「地元産品が他県製品に勝利 満足度7割」の記事を検証する。前8章から、確認すべき項目を6つ選び、それぞれ答えが揃った場合に得られる情報を書け。
\end{quote}

\textbf{解答の指針:} (1)母数。何人が答えたかで7割の重みが変わる。(2)設問。満足度の聞き方と選択肢の数。(3)比較対象。他県製品はどの製品か、誰が選んだか。(4)無回答の扱い。分母から外していれば実力以上の7割かもしれない。(5)調査の主体。地元の利害が関係する調査かどうか。(6)誤差。母数が小さければ差は偶然の範囲かもしれない。6つが埋まれば、この見出しは「どの群を、誰が、どう比べた話か」まで戻せる。

数字は、選んだ人の判断をそのまま映す。だから数字そのものを疑うより、選んだ人に問い掛けるほうが速い。誰が、何と比べて、母数は何か。この3つを、次の記事から唱えてみてほしい。

\end{document}
